% ============================================================================
%  Capitolo 2 -- Primo contatto: assembly RISC-V su PRG32        (Lezione 02)
% ============================================================================
\chapter{Primo contatto: assembly RISC-V su PRG32}
\label{ch:firstcontact}

\epigraphline{Talk is cheap. Show me the code.}
  {Linus Torvalds, mailing list del kernel Linux, 2000: le chiacchiere stanno a zero, fatemi vedere il codice}

\begin{chapgoals}
Alla fine di questo capitolo saprete
\begin{itemize}
\item compilare una cartuccia da un file assembly ed eseguirla su \prg;
\item dire che cos'è un registro, e nominare quelli che si usano ogni giorno;
\item chiamare una funzione del framework, passando gli argomenti nei registri;
\item conservare un valore in memoria da un frame al successivo: caricare, modificare, memorizzare;
\item leggere un'istruzione come il numero di 32 bit che è in realtà;
\item elencare le domande che questo programma solleva, e a cui risponde il resto del libro.
\end{itemize}
\end{chapgoals}

Questo capitolo è il giro «gioca» della spirale per tutta la prima metà del
libro: ogni capitolo successivo rimanda a qualcosa che è stato scritto qui.
\emph{Non} ci si aspetta che il lettore capisca tutto ciò che contiene. Questo
capitolo serve a vedere; i successivi servono a spiegare. Si scrivono quattro
piccoli programmi, ciascuno completo e ciascuno funzionante; sono i passi
verificati del primo laboratorio di assembly del libro di \prg
\citep{prg32textbook}.

% ----------------------------------------------------------------------------
\section{La macchina}

\subsection{I registri, e tre tipi di istruzione}

Un processore \riscv, visto dal programmatore, è una cosa molto piccola.

\begin{description}[style=nextline,leftmargin=1.5em]
\item[Trentadue registri.] Ogni \term{registro} contiene un valore di 32 bit.
  Il processore può calcolare soltanto su valori che stanno nei registri.
\item[Calcolare.] Un'operazione aritmetica o logica tra registri: \ins{add},
  \ins{sub}, \ins{and}, \ins{or}, \ins{xor}, e le loro forme \term{immediate},
  come \ins{addi}, in cui un operando è una piccola costante scritta
  nell'istruzione.
\item[Spostare.] Un trasferimento tra un registro e la memoria: \ins{lw} carica
  una parola, \ins{sw} memorizza una parola. Nient'altro tocca la memoria.
\item[Saltare.] Un cambiamento dell'istruzione che viene eseguita dopo: salti
  condizionati come \ins{beqz}, chiamate e ritorni.
\end{description}

Questa è quasi tutta la macchina. La ricchezza dell'informatica non viene da un
grande insieme di istruzioni, ma dal combinarne uno piccolo in lunghe sequenze.

\subsection{I registri che si usano davvero}

L'hardware chiama i registri da \reg{x0} a \reg{x31} e, con una sola eccezione,
li tratta allo stesso modo. I programmatori li chiamano con altri nomi, che
ricordano il \emph{ruolo} svolto da ciascuno (Tabella~\ref{tab:fc:regs}). I
ruoli non sono fissati dall'hardware ma da un accordo, la \term{convenzione di
chiamata}, parte dell'application binary interface (ABI). Rispettarla è ciò che
permette a una funzione scritta in assembly e al codice C del framework di
chiamarsi a vicenda.

\begin{table}[htb]
\centering\small
\caption{I registri che si usano ogni giorno. La tabella completa è
nell'Appendice~\ref{app:riscv}.}\label{tab:fc:regs}
\begin{tabularx}{\textwidth}{@{}llLl@{}}
\toprule
\textbf{Registro} & \textbf{Nome} & \textbf{Ruolo} & \textbf{Conservato da una chiamata?}\\
\midrule
\reg{a0}--\reg{a7} & argomenti & valori passati a una chiamata; \reg{a0} porta anche il valore di ritorno & no\\
\reg{t0}--\reg{t6} & temporanei & valori di lavoro & no\\
\reg{s0}--\reg{s11} & salvati & valori che una chiamata promette di conservare & sì\\
\reg{ra} & indirizzo di ritorno & dove salta \ins{ret} & no: salvatelo voi\\
\reg{sp} & puntatore allo stack & cima dello stack, allineata a 16 byte & sì\\
\reg{zero} & zero & si legge sempre come 0; le scritture sono ignorate & --\\
\bottomrule
\end{tabularx}
\end{table}

\subsection{Il contratto: tre funzioni}

Un gioco \prg sono tre funzioni che il programmatore scrive e il firmware
chiama (Figura~\ref{fig:fc:contract}). \code{init} viene chiamata una volta,
quando la cartuccia parte, e prepara lo stato del gioco. \code{update} viene
chiamata una volta per frame: legge l'ingresso e cambia lo stato. \code{draw}
viene chiamata una volta per frame: trasforma lo stato in un'immagine. Il
firmware trova le tre funzioni per nome: un prefisso scelto dal programmatore,
seguito da \code{\_init}, \code{\_update} e \code{\_draw}. Questo è l'intero
contratto tra un gioco e la macchina.

\begin{figure}[htb]
\centering
\begin{tikzpicture}[node distance=10mm]
  \node[blkacc,text width=30mm,minimum height=14mm] (i) {\texttt{init}\\[1pt]\footnotesize una volta: prepara lo stato};
  \node[blkamb,text width=30mm,minimum height=14mm,right=of i] (u) {\texttt{update}\\[1pt]\footnotesize a ogni frame: ingresso, stato};
  \node[blkvio,text width=30mm,minimum height=14mm,right=of u] (d) {\texttt{draw}\\[1pt]\footnotesize a ogni frame: dallo stato all'immagine};
  \draw[->,thick] (i) -- (u); \draw[->,thick] (u) -- (d);
  \draw[->,thick,rounded corners=3mm] (d.south) -- ++(0,-6mm) -| node[pos=0.25,below,lbl]{frame successivo} (u.south);
\end{tikzpicture}
\caption{Il contratto tra una cartuccia e il runtime.}
\label{fig:fc:contract}
\end{figure}

% ----------------------------------------------------------------------------
\section{La vostra prima cartuccia}

\subsection{Passo 1: il più piccolo programma che gira}

Il Listato~\ref{lst:fc:step1} è una cartuccia completa. Create una cartella,
scrivete il programma in un file chiamato \fname{step1\_skeleton.S}, e
resistete alla tentazione di copiarlo: imparano anche le dita.

\asmfile[caption={Passo 1: tre funzioni che non fanno nulla, correttamente.},label=lst:fc:step1]{code/l02/c03.S}

Le righe che cominciano con un punto sono \term{direttive}: ordini per
l'assemblatore, non istruzioni per il processore. \code{.option norelax}
mantiene prevedibile la disposizione della cartuccia, e ogni esempio di \prg
comincia così. \code{.section .text} dice che ciò che segue è codice
eseguibile. \code{.global} rende un'\term{etichetta} visibile fuori dal file,
così che il firmware possa trovare le funzioni per nome. Un'etichetta, un nome
seguito da due punti, sta per l'indirizzo di ciò che la segue. L'unica
istruzione del programma, scritta tre volte, è \ins{ret}: ritorna al chiamante.

La cartuccia si compila e si esegue con i comandi incontrati nel
Capitolo~\ref{ch:setup} (Listato~\ref{lst:fc:build}); il file della cartuccia è
identico byte per byte per QEMU, per la scheda e per \prgqt, dove l'indirizzo è
\code{http://localhost:8080}.

\shfile[caption={Compilare ed eseguire una cartuccia.},label=lst:fc:build]{code/l02/c04.sh}

\begin{prgcheck}[Eseguite]
Il risultato atteso è uno schermo nero. Il risultato è reale lo stesso: il
framework ha trovato le tre funzioni, le ha chiamate frame dopo frame, ed è
sopravvissuto.
\end{prgcheck}

\subsection{Passo 2: dire ciao}

Per mettere qualcosa sullo schermo il programma chiama due funzioni del
framework: \code{prg32\_gfx\_clear}, che riempie lo schermo con un colore, e
\code{prg32\_gfx\_text8}, che scrive una stringa con un font di 8 per 8 pixel
(Listato~\ref{lst:fc:step2}, Figura~\ref{fig:fc:hello}). In cima al file due
direttive \code{.equ} danno un nome a due costanti:
\code{.equ BLACK, 0x0000} e \code{.equ WHITE, 0xffff}.

\asmfile[caption={Passo 2: la nuova \code{first\_draw}.},label=lst:fc:step2]{code/l02/c05.S}

\begin{figure}[htb]
\centering
\includegraphics[width=0.55\textwidth]{lab1_step2_hello.png}
\caption{Il passo 2, compilato ed eseguito su \prg.}\label{fig:fc:hello}
\end{figure}

La regola per passare gli argomenti è semplice: il primo argomento va in
\reg{a0}, il secondo in \reg{a1}, e così via. \ins{li} (\emph{load immediate})
mette un numero in un registro; \ins{la} (\emph{load address}) vi mette
l'indirizzo di un'etichetta. La stringa vive nella sezione \code{.rodata},
\emph{read-only data}, dati di sola lettura, e la direttiva \code{.asciz} ne
memorizza i caratteri seguiti da un byte di valore zero che ne segna la fine.

\subsection{Perché salvare \reg{ra}? Lo stack frame}

La prima riga e le ultime tre di \code{first\_draw} sono nuove, e sono il
prezzo da pagare per chiamare un'altra funzione. L'istruzione \ins{call}
memorizza in \reg{ra} l'indirizzo a cui tornare. Ma \code{first\_draw} è stata
a sua volta chiamata dal firmware, e \reg{ra} conteneva già l'indirizzo del
\emph{suo} ritorno: la prima \ins{call} lo sovrascriverebbe, e la \ins{ret}
finale salterebbe nel posto sbagliato.

Una funzione che ne chiama un'altra salva quindi per prima cosa \reg{ra}.
Abbassa \reg{sp} di 16 per fare spazio sullo \term{stack}, una regione di
memoria che ogni funzione può usare per le proprie necessità temporanee, e
memorizza \reg{ra} a \code{12(sp)}, cioè 12 byte sopra l'indirizzo contenuto in
\reg{sp} (Figura~\ref{fig:fc:stack}). Queste righe sono il \term{prologo} della
funzione. Dopo l'ultima chiamata disfa tutto in ordine inverso,
nell'\term{epilogo}: ricarica \reg{ra}, alza \reg{sp} di 16, e ritorna.

\begin{figure}[htb]
\centering
\begin{tikzpicture}[x=1cm,y=0.72cm,font=\ttfamily\small]
  \foreach \i/\t/\o/\f in {0/dati del chiamante/16(sp)/prgrule!60, 1/ra/12(sp)/prgamber!70, 2/libero/{ 8(sp)}/prgpanel, 3/libero/{ 4(sp)}/prgpanel, 4/libero/{ 0(sp)}/prgpanel}{
    \draw[prgsteel,fill=\f] (0,-\i) rectangle (3.6,-\i-1);
    \node at (1.8,-\i-0.5) {\t};
    \node[anchor=west,text=prgcomment] at (3.8,-\i-0.5) {\o};
  }
  \draw[->,very thick,prgred] (-1.2,-4.5) -- (-0.05,-4.5) node[pos=0,left,font=\ttfamily\bfseries] {sp};
  \node[lbl,text=prgcomment] at (1.8,0.4) {indirizzi più alti};
  \node[lbl,text width=7cm,align=center] at (1.8,-5.7) {Lo stack cresce verso il basso: si sottrae da \texttt{sp} per allocare, si somma per rilasciare.};
\end{tikzpicture}
\caption{La memoria dello stack dopo il prologo di \code{first\_draw}.}
\label{fig:fc:stack}
\end{figure}

\begin{prgwarn}[Il classico primo errore]
Dimenticare la coppia \code{sw ra}/\code{lw ra} è il classico primo errore: la
funzione ritorna nel posto sbagliato, e il programma va in crash o si blocca
molto dopo la riga colpevole.
\end{prgwarn}

La convenzione richiede che \reg{sp} resti allineato a un multiplo di 16 byte
tra una chiamata e l'altra, ed è per questo che il frame è di 16 byte anche se
ne vengono usati solo 4. Lo stack è trattato formalmente nei Capitoli
\ref{ch:memory} e~\ref{ch:instructions}.

\subsection{Passo 3: un valore che sopravvive tra i frame}

I registri non conservano i loro valori tra una chiamata di \code{update} e la
successiva: nel frattempo li usa il firmware. Ciò che deve durare deve vivere
in memoria (Listato~\ref{lst:fc:step3}).

\asmfile[caption={Passo 3: la nuova \code{first\_update} e la sua variabile.},label=lst:fc:step3]{code/l02/c06.S}

\code{counter} è un'etichetta: dà un nome a una locazione di memoria, nella
sezione \code{.data}, che contiene una parola di 32 bit, inizialmente zero.
\ins{la} mette in \reg{t0} l'\emph{indirizzo} di quella locazione. \ins{lw}
legge in \reg{t1} il \emph{valore} che si trova a quell'indirizzo; \ins{sw} lo
riscrive. Poiché il processore calcola solo nei registri, ogni modifica di una
variabile ha gli stessi tre movimenti: \emph{caricare, modificare,
memorizzare}. La funzione è una \term{foglia}: non chiama nulla, quindi non ha
bisogno di uno stack frame.

\begin{prgconcept}[Indirizzo o valore? Il punto cruciale della memoria]
Un'etichetta dà un nome a una \emph{locazione}. \ins{la} dà la locazione;
\ins{lw} dà il contenuto (Figura~\ref{fig:fc:addr}). È così che in ultima
analisi funziona ogni variabile di ogni programma. Usare un indirizzo dove
serve un valore, o il contrario, è uno degli errori più comuni dei
principianti: commentate ogni registro come ADDRESS (indirizzo) o VALUE
(valore).
\end{prgconcept}

\begin{figure}[htb]
\centering
\begin{tikzpicture}[font=\ttfamily\small]
  \node[lbl,anchor=west,font=\sffamily\small\bfseries,text=prgsteel] at (0,1.0) {Registri (dentro il processore)};
  \node[anchor=east] at (0.6,0.2) {t0};
  \node[draw=prgaccent,fill=prgaccent!20,minimum width=42mm,minimum height=7mm,anchor=west] (t0) at (0.8,0.2) {indirizzo di counter};
  \node[anchor=east] at (0.6,-1.0) {t1};
  \node[draw=prgamber,fill=prgamber!25,minimum width=42mm,minimum height=7mm,anchor=west] (t1) at (0.8,-1.0) {41, poi 42};
  \node[lbl,anchor=west,font=\sffamily\small\bfseries,text=prgsteel] at (8.2,1.0) {Memoria (sezione \texttt{.data})};
  \node[draw=prgsteel,fill=prgpanel,minimum width=38mm,minimum height=7mm,anchor=west] (m0) at (8.2,0.2) {...};
  \node[draw=prgsteel,fill=prgamber!15,minimum width=38mm,minimum height=7mm,anchor=west] (m1) at (8.2,-0.5) {counter: 41};
  \node[draw=prgsteel,fill=prgpanel,minimum width=38mm,minimum height=7mm,anchor=west] (m2) at (8.2,-1.2) {...};
  \draw[->,very thick,prgaccent] (t0.east) -- node[above,sloped,lbl,text=prgaccent]{\emph{t0 punta qui}} ([yshift=1.5mm]m1.west);
  \draw[->,very thick,prgamber] ([yshift=-1.5mm]m1.west) -- node[below,sloped,lbl,text=prgamber]{\emph{lw copia il contenuto}} (t1.east);
\end{tikzpicture}
\begin{asm}[numbers=none]
la   t0, counter   # t0 = WHERE the value lives
lw   t1, 0(t0)     # t1 = WHAT is stored there
addi t1, t1, 1     # change the copy in the register
sw   t1, 0(t0)     # memory still says 41 until this line
\end{asm}
\caption{Il registro \reg{t0} contiene un indirizzo; \reg{t1} contiene un valore.}
\label{fig:fc:addr}
\end{figure}

Dopo questo passo non cambia nulla di visibile: il contatore avanza, ma nessuno
lo disegna. Come si potrebbe dimostrare che sta avanzando? Con il debugger di
\prgqt, oppure con il passo successivo.

\subsection{Passo 4: rendere visibile il contatore}

Otto righe aggiunte a \code{first\_draw}, dopo la chiamata che scrive il testo,
disegnano una barra verde la cui larghezza è il contatore
(Listato~\ref{lst:fc:step4}, Figura~\ref{fig:fc:bar}). La costante \code{GREEN}
è definita in cima al file con \code{.equ GREEN, 0x07e0}.

\asmfile[caption={Passo 4: le righe aggiunte a \code{first\_draw}.},label=lst:fc:step4]{code/l02/c07.S}

\begin{figure}[htb]
\centering
\includegraphics[width=0.55\textwidth]{lab1_step4_bar.png}
\caption{Passo 4: la barra cresce da 0 a 255 pixel, poi ricomincia.}
\label{fig:fc:bar}
\end{figure}

La funzione \code{prg32\_gfx\_rect} prende cinque argomenti: $x$, $y$,
larghezza, altezza e colore. La larghezza è calcolata da \ins{andi} con 255,
che tiene soltanto gli otto bit bassi del contatore: la barra cresce da 0 a 255
pixel e poi ricomincia. In queste otto righe si nascondono tre capitoli futuri:
l'uguaglianza $255 = \hex{FF}$ (notazione posizionale), \ins{andi} come
maschera di bit (algebra di Boole), e \hex{07E0} come colore (la codifica
dell'informazione).

Il programma completo è il Listato~\ref{lst:fc:full}.

\asmfile[caption={La prima cartuccia completa (\fname{step4\_bar.S}).},label=lst:fc:full]{code/l02/step4_bar.S}

% ----------------------------------------------------------------------------
\section{Leggere ciò che avete scritto}

\subsection{Le istruzioni appena usate}

La Tabella~\ref{tab:fc:ins} ordina per famiglia le istruzioni incontrate
finora. In tutte il risultato va nel \emph{primo} registro scritto. L'insieme
di istruzioni intero di base di \riscv, RV32I, ha una quarantina di istruzioni:
un quarto è già stato usato.

\begin{table}[htb]
\centering\small
\caption{Le istruzioni usate in questo capitolo.}\label{tab:fc:ins}
\begin{tabularx}{\textwidth}{@{}llLl@{}}
\toprule
\textbf{Famiglia} & \textbf{Istruzione} & \textbf{Che cosa fa} & \textbf{Esempio}\\
\midrule
Aritmetica & \ins{add}, \ins{sub} & combina due registri & \code{add t0, t1, t2}\\
Logica & \ins{and}, \ins{or}, \ins{xor} & operazione bit a bit su due registri & \code{and t0, t1, t2}\\
Immediato & \ins{addi}, \ins{andi}, \ins{ori}, \ins{xori} & un registro e una piccola costante & \code{addi t1, t1, 1}\\
Memoria & \ins{lw}, \ins{sw} & carica o memorizza una parola di 32 bit & \code{lw t1, 0(t0)}\\
Controllo & \ins{call}, \ins{ret} & va a una funzione e ne ritorna & \code{call prg32\_gfx\_rect}\\
\bottomrule
\end{tabularx}
\end{table}

\subsection{Pseudo-istruzioni}

Parecchi degli mnemonici usati non sono affatto istruzioni del processore. Sono
\term{pseudo-istruzioni}: scorciatoie che l'assemblatore espande in una o due
istruzioni vere (Tabella~\ref{tab:fc:pseudo}). Una pseudo-istruzione non è
un'operazione hardware distinta. È il registro \reg{zero} a renderne possibili
molte: un piccolo insieme di istruzioni, riusato con cura.

\begin{table}[htb]
\centering\small
\caption{Le pseudo-istruzioni e le loro espansioni.}\label{tab:fc:pseudo}
\begin{tabular}{@{}lll@{}}
\toprule
\textbf{Voi scrivete} & \textbf{L'assemblatore emette} & \textbf{Significato}\\
\midrule
\code{li t0, 42} & \code{addi t0, zero, 42} & carica una costante\\
\code{mv t0, t1} & \code{addi t0, t1, 0} & copia un registro\\
\code{not t0, t1} & \code{xori t0, t1, -1} & complemento bit a bit\\
\code{la t0, counter} & \code{auipc t0, \ldots} + \code{addi t0, t0, \ldots} & carica l'indirizzo di un'etichetta\\
\code{ret} & \code{jalr zero, 0(ra)} & salta all'indirizzo contenuto in \reg{ra}\\
\code{beqz t0, skip} & \code{beq t0, zero, skip} & salta se \reg{t0} vale zero\\
\bottomrule
\end{tabular}
\end{table}

\subsection{Un'istruzione è un numero di 32 bit}

Il processore non legge testo. L'assemblatore traduce ogni riga in una parola
di 32 bit, ed è quella parola a essere memorizzata nella cartuccia ed eseguita.
La Figura~\ref{fig:fc:encoding} seziona \code{addi t1, t1, 1}: dodici bit
contengono la costante, cinque bit indicano il registro sorgente (\reg{t1} è
\reg{x6}), tre bit selezionano l'operazione, cinque bit indicano il registro di
destinazione e sette bit, l'\term{opcode}, dicono di che tipo di istruzione si
tratta. Scritti in esadecimale, i trentadue bit sono \hex{00130313}.

\begin{figure}[htb]
\centering
\begin{bytefield}[bitwidth=0.029\textwidth,bitheight=8mm]{32}
\bitheader[endianness=big]{0,6,7,11,12,14,15,19,20,31}\\
\bitbox{12}[bgcolor=prgamber!30]{\texttt{000000000001}} & \bitbox{5}[bgcolor=prgaccent!20]{\texttt{00110}} & \bitbox{3}[bgcolor=prgviolet!20]{\texttt{000}} & \bitbox{5}[bgcolor=prggreen!20]{\texttt{00110}} & \bitbox{7}[bgcolor=prgsteel!20]{\texttt{0010011}}\\
\bitbox[]{12}{\footnotesize immediato\\la costante 1} & \bitbox[]{5}{\footnotesize rs1\\\texttt{t1}=\texttt{x6}} & \bitbox[]{3}{\footnotesize funct3\\add} & \bitbox[]{5}{\footnotesize rd\\\texttt{t1}=\texttt{x6}} & \bitbox[]{7}{\footnotesize opcode\\reg.-imm.}
\end{bytefield}
\caption{\code{addi t1, t1, 1} è il numero \hex{00130313}.}
\label{fig:fc:encoding}
\end{figure}

\begin{prgconcept}
L'assembly è una scrittura leggibile del codice macchina. Un'istruzione è un
dato: il processore ne legge i campi con circuiti logici, che verranno
progettati nei capitoli sulla logica e sul processore.
\end{prgconcept}

\subsection{PRG32-QT: un microscopio per i registri}

\prgqt esegue la stessa cartuccia sul desktop un'istruzione alla volta, si
ferma tra due istruzioni qualunque e mostra ogni registro.

\begin{prgtry}[Eseguite passo-passo \code{first\_update}]
Caricate il passo~3 in modalità debug, mettete un breakpoint su
\code{first\_update} e premete \key{Step} quattro volte. Prima di ogni
pressione, dite che cosa conterranno \reg{t0} e \reg{t1}; dopo, verificate.
\end{prgtry}

\subsection{Tracciare a mano}

Lo stesso esercizio si può, e si deve, fare su carta. La
Tabella~\ref{tab:fc:trace} segue \code{first\_update} per alcuni frame. Al
frame 256 il contatore continua a crescere, ma i suoi otto bit bassi tornano a
zero, perché $256 = 1\,0000\,0000_2$: la barra ricomincia. Tracciare a mano è
la più preziosa abilità di debug insegnata in questo libro.

\begin{table}[htb]
\centering\small
\caption{Tre frame, e uno in più, di \code{first\_update}.}\label{tab:fc:trace}
\begin{tabular}{@{}ccccc@{}}
\toprule
\textbf{Frame} & \textbf{\ins{lw}: \reg{t1} diventa} & \textbf{\ins{addi}: \reg{t1} diventa} & \textbf{\ins{sw}: counter diventa} & \textbf{Larghezza}\\
\midrule
1 & 0 & 1 & 1 & 1\\
2 & 1 & 2 & 2 & 2\\
256 & 255 & 256 & 256 & 0\\
257 & 256 & 257 & 257 & 1\\
\bottomrule
\end{tabular}
\end{table}

% ----------------------------------------------------------------------------
\section{Che cosa viene dopo}

Il programma di questo capitolo ha sollevato più domande di quante ne abbia
risolte. La Tabella~\ref{tab:fc:questions} è la mappa in avanti del libro: ogni
capitolo successivo si apre indicandone una riga.

\begin{table}[htb]
\centering\small
\caption{Le domande che il primo programma ha già sollevato.}\label{tab:fc:questions}
\begin{tabularx}{\textwidth}{@{}P{0.29\textwidth}LP{0.2\textwidth}@{}}
\toprule
\textbf{Avete visto} & \textbf{La domanda} & \textbf{Risposta in}\\
\midrule
\hex{07e0}, \hex{ffff}, 255 & Perché i programmatori scrivono i numeri in esadecimale? & Capitoli \ref{ch:information}, \ref{ch:integers}\\
una posizione sotto lo zero mostrata come \hex{FFFFFFFF} & Come fa un registro a contenere un numero negativo? & Capitolo \ref{ch:integers}\\
una palla che si muove di meno di un pixel per frame & Come si rappresentano le frazioni senza un'unità in virgola mobile? & Capitolo \ref{ch:reals}\\
\code{andi a2, t1, 255} & Che cosa fa AND a trentadue bit tutti insieme? & Capitoli \ref{ch:boole1}, \ref{ch:boole2}\\
\code{addi t1, t1, 1} & Quale circuito somma davvero due numeri? & Capitolo \ref{ch:boole2}\\
un pulsante premuto è un 1 & Di che cosa è fatto, fisicamente, un bit? & Capitolo \ref{ch:electronics}\\
\code{counter} sopravvive tra i frame & Come fa un circuito a ricordare un bit? & Capitoli \ref{ch:sequential}, \ref{ch:memory}\\
\hex{00130313} & Come fa il processore a decodificarla ed eseguirla? & Capitoli \ref{ch:processor}, \ref{ch:instructions}\\
tre funzioni e un rettangolo & Se ne può fare un gioco intero? & \vermone: Capitoli \ref{ch:verm1p1}, \ref{ch:verm1p2}, \ref{ch:verm1p3}, \ref{ch:verm1p4}\\
\bottomrule
\end{tabularx}
\end{table}

% ----------------------------------------------------------------------------
% Generated by tools/textbook/export_listings.js from decks/plan.js and code/l02.S: do not edit.
\section{Il programma di questo capitolo}\label{sec:lab:l02}

I listati di questo capitolo sono parti di un solo programma,
\fname{code/l02.S} nel repository del corso, stampato per intero nel
Listato~\ref{lst:lab:l02}. I test automatici del materiale del corso lo
compilano e lo eseguono sull'esecutore senza interfaccia di \prgqt{}, con
ingressi programmati, e controllano il frame; la
Figura~\ref{fig:lab:l02} è il frame scritto dal test.

\begin{description}[style=nextline,leftmargin=1.5em]
\item[Sullo schermo.] HELLO PRG32 e, sotto, una barra verde che cresce di un pixel per frame e ricomincia a 256.
\item[Che cosa premere.] Niente: il programma va da solo.
\end{description}

\begin{figure}[htb]
\centering
\includegraphics[width=0.5\textwidth]{lab_l02.png}
\caption{\fname{l02.S} su \prgqt{}.}\label{fig:lab:l02}
\end{figure}

\asmfile[caption={Il programma completo del capitolo, \fname{code/l02.S}.},label=lst:lab:l02]{code/l02/program.S}

Per eseguirlo su \prgqt{} i comandi sono quelli del
Listato~\ref{lst:lab:l02:run}; il Paragrafo~\ref{sec:prg:testing} dà gli
stessi passi per QEMU e per la scheda.

\begin{shell}[caption={Compilare \fname{l02.S} ed eseguirlo su \prgqt{}.},label=lst:lab:l02:run]
# 1. start PRG32-QT on this computer (it listens on port 8080)

# 2. in the PRG32 folder, with ESP-IDF active: build the cartridge
python3 -m prg32 cartridge build ~/arc12/code/l02.S --portable \
    --entry-prefix first --name l02 --out l02.prg32

# 3. send it to PRG32-QT and start it
python3 -m prg32 esp32c6 upload-and-run l02.prg32 \
    --url http://localhost:8080 --slot cart0

# 4. look, press, compare with the previous slide
\end{shell}


% ----------------------------------------------------------------------------
\section{Verificate di aver capito}

\begin{quiz}{Dopo \code{la t0, counter} e \code{lw t1, 0(t0)}, che cosa contiene \reg{t0}?}
\item Il valore del contatore
\item L'indirizzo del contatore
\item Zero, perché \ins{la} azzera il registro
\item L'indirizzo di ritorno
\end{quiz}

\begin{quiz}{Una funzione che chiama \code{prg32\_gfx\_clear} deve salvare un registro prima della chiamata. Quale?}
\item \reg{a0}, perché porta l'argomento
\item \reg{sp}, perché lo stack si sposta
\item \reg{ra}, perché \ins{call} lo sovrascrive
\item \reg{t0}, perché è temporaneo
\end{quiz}

\begin{quiz}{\reg{t1} contiene 300. Che cosa lascia in \reg{a2} l'istruzione \code{andi a2, t1, 255}?}
\item 300
\item 255
\item 45
\item 44
\end{quiz}

\section{Esercizi}

\begin{exercise}[su \prg]
\begin{enumerate}[nosep]
\item Cambiate la stringa del titolo e il suo colore in giallo (\hex{ffe0});
  eseguite di nuovo.
\item Fate crescere la barra a velocità doppia, poi quadrupla, cambiando una
  sola istruzione.
\item Fate accorciare la barra invece di farla crescere. Che cosa succede
  quando il contatore passa per lo zero?
\item In \prgqt eseguite passo-passo \code{first\_update} e annotate \reg{t0} e
  \reg{t1} dopo ogni istruzione.
\item Disegnate una seconda barra, rossa, la cui larghezza è il contatore
  modulo 64.
\end{enumerate}
\expected{Nel terzo compito il contatore passa da 0 a \hex{FFFFFFFF}: con
\code{andi 255} la barra salta a 255 e poi si accorcia. Conservate
l'osservazione: la spiega il Capitolo~\ref{ch:integers}.}
\end{exercise}

\section{Glossario}

\begin{glossarytable}
register & registro\\
load / store & caricare / memorizzare\\
address & indirizzo\\
immediate & immediato (costante)\\
branch & salto condizionato\\
label & etichetta\\
stack & pila (stack)\\
stack pointer & puntatore alla pila\\
return address & indirizzo di ritorno\\
calling convention & convenzione di chiamata\\
machine code & codice macchina\\
cartridge & cartuccia\\
\end{glossarytable}

\section{Riepilogo e letture}

Un processore \riscv sono trentadue registri e un piccolo insieme di istruzioni
che calcolano, spostano e saltano. Un gioco \prg sono tre funzioni; gli
argomenti viaggiano nei registri da \reg{a0} ad \reg{a7}; una funzione che ne
chiama un'altra salva \reg{ra}. Le variabili vivono in memoria: \ins{la} per
l'indirizzo, \ins{lw} e \ins{sw} per il valore. Un'istruzione è un numero di 32
bit; l'assembly ne è la scrittura leggibile.

\paragraph{Letture.} \emph{Playing with RISC-V}, capitolo~5, «First Contact»
\citep{prg32textbook}; \citet{patterson2020}, capitolo~2; il tutorial
\fname{docs/learn/tutorial\_ascii\_game.md} del repository di \prg.